\begin{textsample}{POS Dim 3 – se – Score 18.00 – se/se\_inf\_5.txt}  \label{ex:f3_pos_163}
2.3.4 Professores na Educação \textbf{Infantil} : Qualidade da Docência Ter clareza da importância dos professores na garantia dos direitos das \textbf{crianças} a partir de uma concepção de \textbf{infância}, bem como de Educação \textbf{Infantil}, é ponto de partida para a construção de um trabalho pedagógico consistente, balizado pela perspectiva de que a docência na Educação \textbf{Infantil} tem papel fundamental na promoção da \textbf{criança} ao desenvolvimento de sua humanidade a partir do incentivo à autonomia \textbf{infantil} nos diversos tempos e espaços educativos.
O ato de \textbf{cuidar} tão presente nas ações docentes na Educação \textbf{infantil} toma forma para além do simples limpar, alimentar... “ \textbf{cuidar} significa também ensinar, produzir o humano no próprio corpo da \textbf{criança} e sua relação com ele, passando pela alimentação, pelo andar, movimentar -se, etc ” ( ARCE, 2013, p.33 ), ou seja, o professor(a) promove na \textbf{criança} a sua segunda natureza com o \textbf{nascimento} para o mundo social.
Os conhecimentos sobre as \textbf{crianças} são fundamentais ao professor para o estabelecimento das proposições curriculares na Educação \textbf{Infantil} com a valorização da \textbf{infância} e, sobretudo, com um projeto educativo voltado à formação das novas gerações de cada comunidade, exercendo o papel social da \textbf{instituição} educacional.
Nesse sentido, > Cabe à \textbf{professora} e ao professor criar oportunidades para que a \textbf{criança} no processo de elaborar sentidos pessoais, se aproprie de elementos significativos de sua cultura não como verdades absolutas, mas como elaborações dinâmicas e provisórias.
Trabalha -se com os saberes da prática que as \textbf{crianças} vão construindo ao mesmo tempo em que se garante a apropriação ou a construção por ela de novos conhecimentos.
Para tanto a \textbf{professora} e o professor observará as ações \textbf{infantis}, individuais e coletivas, \textbf{acolhem} suas perguntas e suas respostas, buscam compreender o significado de sua conduta. ( BRASIL, CNE/CEB, 2009 ).
O professor e a \textbf{professora} têm um papel fundamental, tanto na investigação dos processos de significação das \textbf{crianças}, quanto na escolha de atividades promotoras de desenvolvimento.
Ele deve ser responsável por criar bons contextos de mediação entre \textbf{crianças}, seu entorno social e os vários elementos da cultura, cabe -lhes a arte e a competência de criar condições para que as aprendizagens ocorram tanto nas \textbf{brincadeiras} livres quanto nas demais situações que as próprias \textbf{crianças} estabelecem enquanto \textbf{brincam}, produzem e aprendem cooperativamente.
Essa tarefa docente implica uma forma de compreender a relação professores e alunos a partir da superação de uma relação verticalizada de subordinação e dominação, passando a reconhecer a trilha do diálogo em estreita relação de alteridade entre \textbf{adultos} e a \textbf{infância} carente de sentidos e significados sobre os fenômenos que precisam se apropriar ( BUSS-SIMÃO e ROCHA, 2017, p.
É o professor quem planeja as melhores atividades, aproveita as diversas situações do cotidiano e potencializa as interações.
Apresentando às \textbf{crianças} o mundo em sua complexidade : a natureza, a sociedade, os \textbf{sons}, os jogos, as \textbf{brincadeiras}, os conhecimentos construídos ao longo da sua vida, dando condições para a construção de sua identidade, autonomia dentro de um grupo social.
No trabalho docente a \textbf{criança} é elevada a \textbf{centralidade} das relações de sociabilidade, tornando -se imperativo que uma efetiva ação educativa se constitua na realização de importantes elementos para o trabalho pedagógico, como : observação permanente e sistemática, o \textbf{registro} e a documentação, como forma de avaliar o proposto, conhecer o vivido e repropor as experiências a serem privilegiadas e as formas de organização dos espaços, dos tempos e dos materiais para dar conta dos princípios que norteiam o desenvolvimento e a educação das \textbf{crianças} até seis anos de idade em \textbf{instituições} educativas, quais sejam : as interações, as \textbf{brincadeiras} e as linguagens.
Outra importante tarefa das \textbf{instituições} de Educação \textbf{Infantil} é o envolvimento da família dos alunos na \textbf{rotina} da unidade escolar que é determinante no processo de aprendizagem da \textbf{criança} e se torna a base para a promoção do desenvolvimento das \textbf{crianças} em sua integralidade.
É possível integrar o conhecimento das famílias nos projetos e demais atividades pedagógicas.
Não só as questões culturais e regionais podem ser inseridas nas programações por meio da participação dos pais e demais familiares, mas também as questões afetivas e motivações familiares devem fazer parte do cotidiano pedagógico.
Para Kramer ( 2006 ) o trabalho com os \textbf{bebês} e \textbf{crianças} \textbf{pequenas} nos espaços da Educação \textbf{Infantil} implica em uma multiplicidade de processos educativos e uma formação profissional que evidencie a relevância de se discutir quais elementos devem ser contemplados nesse processo de formação.
A formação do professor de Educação \textbf{Infantil} ( LDB/1996 ) consolida o \textbf{cuidar} / educar \textbf{crianças} \textbf{pequenas} através de uma perspectiva ampla e indissociável. \textbf{Cuidar} e educar são ações intrínsecas e de responsabilidade da família, da escola, e consequentemente, dos professores.
Para Vital Didonet “ Todos têm de saber que só se \textbf{cuida} educando e só se educa \textbf{cuidando} ”.
A formação de professoras(es) da Educação \textbf{Infantil} é tarefa recente na história da educação brasileira, este lugar do professor da educação \textbf{infantil} é desafiador na perspectiva de se pensar uma prática pedagógica que promova nas \textbf{crianças} a ocupação de um lugar importante na formação do ser humano.
A formação desses profissionais da Educação \textbf{Infantil} deve partir do reconhecimento da \textbf{criança} como um sujeito histórico e de direitos, que deve acessar o conhecimento científico imbricado pelo diálogo com o cotidiano do qual faz parte.
Segundo Corsaro ( 2011, p. 63 ), em sua teoria sobre a “ interpretação reprodutiva das \textbf{crianças} ”, elas \textbf{assimilam}, reproduzem e interpretam o cotidiano dos \textbf{adultos}.
Nesse sentido, elas também constroem e modificam seus espaços, sendo por isso, consideradas sujeitos de direitos e de cultura.

% matched lemmas: acolher, adulto, assimilar, bebê, brincadeira, brincar, centralidade, criança, cuidar, infantil, infância, instituição, nascimento, pequeno, professora, registro, rotina, som
\end{textsample}
